\chapter{Short Interest, Borrow and Crowding}\label{ch:s1:short-interest-borrow-and-crowding}

In January 2021 the short interest in GameStop reached 122.97\% of its float: more shares had been sold short than were available to trade, because the
same shares had been lent, sold and lent again. On January 27 the stock closed at \$347.51, more than 1,600\% above its close sixteen days earlier. The SEC's
staff later found that short sellers covering their positions had contributed to some of the sharpest rises, though their buying was a small fraction of
all buying. Short interest is usually read as a signal, since short sellers tend to be informed; in a crowded name it is also a position that can be forced
shut. This chapter builds a lending market on the synthetic stocks, measures what short interest and borrow fees predict, what the fees cost, and what a
crowded short book can lose. The build is \texttt{firm.lendsig}.

\section{Short interest as a signal}

Short interest (Book~1, chapter~16) is the number of shares sold short and not yet bought back, reported as a share of the shares outstanding or of the float,
or in days of average volume (days to cover). The evidence that it predicts returns is strong in some forms and fragile in others. Boehmer, Jones and Zhang,
with daily NYSE short-sale data for 2000--2004, found that heavily shorted stocks underperformed lightly shorted ones by a risk-adjusted 1.16\% over the next
twenty trading days: short sellers as a group are informed. Asquith, Pathak and Ritter, with monthly short interest over a longer period, found many of the
documented patterns not robust: equally weighted portfolios of high-short-interest stocks generally underperformed, value-weighted ones did not.

\begin{dated}{2026-09}{Short-interest reporting in the United States}\label{dat:s1:short-interest-borrow-and-crowding:finra}
Under FINRA Rule 4560, broker-dealers report the gross short positions in all equity securities held in their customer and proprietary accounts twice a
month, as of settlement dates FINRA designates; reports are due by 6~p.m.\ Eastern Time on the second business day after the settlement date, and FINRA
publishes the aggregated short interest on a published schedule some days later. A short-interest signal therefore sees positions as they stood one to
several weeks earlier; daily short-sale volume and securities-lending data are the faster, commercial alternatives.
\end{dated}

The synthetic market has no short sellers, so \texttt{firm.lendsig} adds them. Each stock has a lendable supply, a lognormal share of its shares centred on
20\%. Short demand has two parts: informed shorts, who follow the stock's planted expected return with a twenty-day lag and short in proportion to how
negative it is; and uninformed shorts, persistent noise with a median of 2\% of shares. Short interest is the demand the supply can meet. Over years 3 to
10 the median stock has 3.1\% of its shares short, the 90th percentile 7.5\% and the most shorted 54\%; median days to cover is 2.1, and the 90th
percentile 6.8. Because informed shorts are in the demand by construction, short interest predicts returns: minus short interest has a monthly rank IC of
0.061, days to cover 0.048, and a crowding score (the average rank of short interest, utilisation and days to cover) 0.058.

\section{Borrow fees and utilisation}

\begin{definition}[Borrow-fee signal]\label{def:s1:short-interest-borrow-and-crowding:fee}
A \emph{borrow-fee signal}\index{borrow-fee signal} ranks stocks on the fee to borrow them for a short sale (or on utilisation, the share of lendable
supply on loan); high fees mark stocks in heavy short demand relative to supply, which have tended to underperform.
\end{definition}

The synthetic fee is 0.25\% a year (general collateral) until utilisation reaches 60\%, and then rises steeply to 40\% at full utilisation. Only 2.9\% of
stock-days pay more than 1\%, but the top percentile pays the maximum. The fee is the cost of the signal, and here it decides the result:

\begin{center}
\small
\begin{tabular}{@{}lccc@{}}
\toprule
decile books, rebuilt monthly & short interest & crowding & borrow fee\\
\midrule
rank IC against the next month & 0.061 & 0.058 & 0.006\\
Sharpe ratio before fees & 3.33 & 3.30 & 1.18\\
Sharpe ratio after fees & 2.24 & 2.08 & $-0.66$\\
Sharpe ratio after fees and costs & 1.84 & 1.69 & $-0.82$\\
gross return a year & 8.3\% & 8.2\% & 2.5\%\\
borrow fees; trading costs (a year) & 2.7\%; 0.95\% & 3.0\%; 0.94\% & 4.0\%; 0.33\%\\
average fee on the short side & 5.4\% & 6.1\% & 7.9\%\\
\bottomrule
\end{tabular}
\end{center}

The short-interest book keeps most of its edge after fees (8.3\% gross, 4.6\% net of fees and costs). The fee book is the warning: the stocks with the
highest fees do underperform, but by less than it costs to short them, and the book loses. Engelberg, Reed and Ringgenberg made the related point on real data
that the fee's level is not the only risk: stocks whose loans can become expensive or be recalled carry short-selling risk, and have lower returns, less
price efficiency and less short selling.

\begin{omfigure}
\begin{tikzpicture}
\begin{axis}[width=10cm,height=5.2cm,xlabel={\small year},ylabel={\small cumulative return (\%)},tick label style={font=\footnotesize},
  xmin=2,xmax=10,grid=major,grid style={gray!25},legend pos=north west,legend style={font=\scriptsize},table/col sep=comma]
\addplot[black!50,thick,dashed] table[x=year,y=gross]{figdata/strategies-1/08-short-interest-borrow-and-crowding/book.csv};
\addplot[omDef!60,thick] table[x=year,y=after_fee]{figdata/strategies-1/08-short-interest-borrow-and-crowding/book.csv};
\addplot[omDef,very thick] table[x=year,y=net]{figdata/strategies-1/08-short-interest-borrow-and-crowding/book.csv};
\legend{before borrow fees,after borrow fees,after fees and trading costs}
\end{axis}
\end{tikzpicture}

\omcaption{The short-interest decile book on the synthetic lending market: cumulative return before borrow fees, after them, and after trading costs too
(gross exposure one). Data: \texttt{s1\_lending.book}.}\label{fig:s1:short-interest-borrow-and-crowding:book}
\end{omfigure}

The safer use of the signal is avoidance. A long-only equal-weighted portfolio of the synthetic stocks earns 9.1\% a year; the same portfolio without its most
shorted decile earns 10.0\%, a difference of 0.94\% a year with a $t$ statistic of 7.2, and it pays no fees at all.

\section{Recalls, squeezes and the lending market's risk}

\begin{definition}[Squeeze risk]\label{def:s1:short-interest-borrow-and-crowding:squeeze}
\emph{Squeeze risk}\index{squeeze risk} is the risk that short sellers of a stock are forced to buy back at the same time (because loans are recalled,
margin is called or losses mount), pushing the price up and forcing further buying.
\end{definition}

A short position has three risks a long position does not: the fee can rise, the loan can be recalled, and the loss is unbounded. In the synthetic market,
loans are recalled with a daily probability of 0.05\% at low utilisation, rising to 2\% at full utilisation: the names a short-interest signal points to are
the names most likely to be taken away. The squeeze is the extreme case. The SEC's staff described its mechanics for GameStop: margin calls or losses make
short sellers buy, their buying raises the price, and the rise forces more short sellers to buy; they also found the price stayed high after the covering
had waned, so that the squeeze explained part of the episode, not all of it.

\section{Crowded shorts}

\begin{definition}[Short-crowding measure]\label{def:s1:short-interest-borrow-and-crowding:crowding}
A \emph{short-crowding measure}\index{short-crowding measure} scores how many short sellers are in a stock relative to what they can get out through,
combining short interest, utilisation of lendable supply and days to cover.
\end{definition}

The squeeze stress (\texttt{firm.lendsig.squeeze\_loss}) takes a book short the twenty most crowded
names, half its capital in all, at each month-end, and asks what it loses if half of each name's short interest is bought back over five days, with the
price following the square-root impact law (Book~7, chapter~27) for as long as the buying lasts. The median loss is 4.0\% of capital and the worst 8.3\%;
the same book's worst actual five-day loss over the eight years was 3.0\%. If only a quarter is covered the median is 2.8\%; if half is covered over ten
days, 5.7\% and at worst 12.3\%. The square-root law was not built for buying many days of volume, so these are floors, not forecasts: the GameStop price
rose about 2,700\% from its low to its high. At that scale the arithmetic is simpler than any model: one name rising twentyfold in a book short twenty names
equally with half its capital costs 47.5\% of capital, from one name.

\begin{omfigure}
\begin{tikzpicture}
\begin{axis}[ybar,bar width=6pt,width=11cm,height=5.0cm,xlabel={\small five-day return of the crowded short book (\% of capital)},
  ylabel={\small month-ends},tick label style={font=\footnotesize},xmin=-10,xmax=3.5,ymin=0,grid=major,grid style={gray!25},
  legend style={font=\scriptsize,at={(0.5,1.03)},anchor=south,legend columns=2},table/col sep=comma]
\addplot[fill=omThm!60,draw=omThm] table[x=bin,y=squeeze]{figdata/strategies-1/08-short-interest-borrow-and-crowding/squeeze.csv};
\addplot[fill=omDef!60,draw=omDef] table[x=bin,y=actual]{figdata/strategies-1/08-short-interest-borrow-and-crowding/squeeze.csv};
\legend{if half the short interest covers in five days,actual next five days}
\end{axis}
\end{tikzpicture}

\omcaption{The book short the twenty most crowded synthetic names ($-0.5$ of capital): its actual return over the five days after each month-end, and its loss
in the squeeze stress, over 96 month-ends. Data: \texttt{s1\_lending.squeeze}.}\label{fig:s1:short-interest-borrow-and-crowding:squeeze}
\end{omfigure}

\section{Strategy files}

\begin{strategyfile}{High short interest avoidance}\label{strat:s1:short-interest-borrow-and-crowding:avoid}
\sfield{Who pays you, and why} Informed short sellers, whose positions reveal information; the long-only investor free-rides on it.
\sfield{Instruments and venues} Stocks, long only.
\sfield{Signal} Short interest as a share of shares or float, as last published.
\sfield{Sizing and execution} Exclude or underweight the most-shorted names in an existing portfolio.
\sfield{Costs} None of its own beyond the rebalancing.
\sfield{How it dies} When short interest stops being informative; value-weighted evidence is weak.
\sfield{Horizon, capacity, infrastructure} Months; very large capacity; the published short-interest file.
\sfield{Backtest honestly} Short interest as of its publication date, not its settlement date; equal and value weights both.
\sfield{Sources} Asquith, Pathak and Ritter (2004/2005); this chapter's simulation (0.94\% a year, $t = 7.2$, a planted effect).
\end{strategyfile}

\begin{strategyfile}{Borrow-fee signal}\label{strat:s1:short-interest-borrow-and-crowding:fee}
\sfield{Who pays you, and why} Short demand exceeding supply marks overpricing that shorts cannot fully correct.
\sfield{Instruments and venues} Stocks, with a securities-lending feed.
\sfield{Signal} The borrow fee or utilisation, daily.
\sfield{Sizing and execution} As an avoidance or underweight signal; shorting the names themselves pays the fee.
\sfield{Costs} The fee itself: on the synthetic market it exceeds the underperformance.
\sfield{How it dies} On the short side, by construction; recalls and fee spikes.
\sfield{Horizon, capacity, infrastructure} Weeks to months; commercial lending data.
\sfield{Backtest honestly} Fees charged on every short day at the rate then prevailing; recalls simulated.
\sfield{Sources} Engelberg, Reed and Ringgenberg (2018); this chapter's simulation.
\end{strategyfile}

\begin{strategyfile}{Days-to-cover ranking}\label{strat:s1:short-interest-borrow-and-crowding:dtc}
\sfield{Who pays you, and why} As for short interest, weighted by how hard the position is to exit.
\sfield{Instruments and venues} Stocks.
\sfield{Signal} Short interest divided by average daily volume.
\sfield{Sizing and execution} Long low, short high, or avoidance.
\sfield{Costs} Borrow fees on the short side are highest here.
\sfield{How it dies} Squeezes in the names it shorts.
\sfield{Horizon, capacity, infrastructure} Weeks to months.
\sfield{Backtest honestly} Volume and short interest from the same dates; fees and recalls.
\sfield{Sources} Boehmer, Jones and Zhang (2008): heavily shorted stocks underperformed by a risk-adjusted 1.16\% over 20 days, 2000--2004.
\end{strategyfile}

\begin{strategyfile}{Crowded-short risk overlay}\label{strat:s1:short-interest-borrow-and-crowding:overlay}
\sfield{Who pays you, and why} Nobody: it is insurance, paid for by giving up part of the short-interest premium.
\sfield{Instruments and venues} The firm's short books.
\sfield{Signal} The crowding score of each short position.
\sfield{Sizing and execution} Cap the book's exposure to crowded shorts; limit each name by days to cover; buy calls on the most crowded names.
\sfield{Costs} The premium forgone and any options bought.
\sfield{How it dies} It does not; it costs.
\sfield{Horizon, capacity, infrastructure} Daily monitoring; lending data.
\sfield{Backtest honestly} Stress with squeeze scenarios, not only with history, which rarely contains the tail.
\sfield{Sources} SEC staff report (2021) on GameStop; this chapter's squeeze stress.
\end{strategyfile}

\section{Tutorial: January 2021}\label{tut:s1:short-interest-borrow-and-crowding:tut}

\begin{tutorial}
\textbf{Goal.} Add a lending market to the synthetic stocks, measure the short-interest and fee signals before and after fees, and stress a crowded short
book. \textbf{End state:} the table, \cref{fig:s1:short-interest-borrow-and-crowding:book} and \cref{fig:s1:short-interest-borrow-and-crowding:squeeze}.
\begin{tutsteps}
\item \textbf{The lending market}: supply, informed and uninformed demand, fees and recalls.
\omcode{code/firm/lendsig/firm_lendsig.py}{48}{75}{The lending layer.}\label{lst:s1:short-interest-borrow-and-crowding:lending}
\item \textbf{The book}: a decile book with fees on its shorts.
\omcode{code/strategies-1/08-short-interest-borrow-and-crowding/python/s1_lending.py}{86}{94}{Borrow fees charged on the short side.}\label{lst:s1:short-interest-borrow-and-crowding:book}
\item \textbf{Run} \texttt{levels}, \texttt{ic}, \texttt{book}, \texttt{avoidance}, \texttt{squeeze} and \texttt{fig\_lending.py}.
\end{tutsteps}
\textbf{What to change next.} Make recalls close positions and charge their re-entry; let uninformed shorts crowd into the informed names (herding); price
the squeeze with a feedback term in which covering raises the price and the price raises covering.
\end{tutorial}

\section{Build: the lending market}\label{bld:s1:short-interest-borrow-and-crowding:lendsig}

\begin{build}
\bfield{Purpose} A synthetic stock-lending market and the short-selling signals and stresses that need one.
\bfield{Interface} \texttt{LendingConfig}, \texttt{simulate\_lending(alpha, listed, shares, cfg)}, \texttt{fee\_of(util, cfg)}, \texttt{days\_to\_cover(si,
shares, adv\_shares)}, \texttt{crowding(si, util, dtc)}, \texttt{squeeze\_loss(w, si, shares, adv\_shares, sigma, cover, days, eta)}.
\bfield{Rules} Short interest never above supply; fees charged daily on shorts; recalls random with a probability rising in utilisation.
\bfield{Acceptance tests} \texttt{code/firm/lendsig/tests/}: the fee schedule by hand; informed demand concentrated in negative-alpha names; days to cover,
the crowding score and the squeeze loss by hand.
\bfield{Stretch} Rehypothecation (short interest above 100\% of float); lending supply from index funds; a feedback squeeze.
\end{build}

\begin{omsources}
\item US SEC staff, \emph{Staff Report on Equity and Options Market Structure Conditions in Early 2021}, October 2021.
\item P.~Asquith, P.~A.~Pathak and J.~R.~Ritter, ``Short interest, institutional ownership, and stock returns'', \emph{Journal of Financial Economics}
78(2), 2005.
\item E.~Boehmer, C.~M.~Jones and X.~Zhang, ``Which shorts are informed?'', \emph{Journal of Finance} 63(2), 2008.
\item J.~E.~Engelberg, A.~V.~Reed and M.~C.~Ringgenberg, ``Short-selling risk'', \emph{Journal of Finance} 73(2), 2018.
\item FINRA Rule 4560, short-interest reporting.
\end{omsources}

\section{Exercises}

\begin{exercise}[$\star$]\label{exo:s1:short-interest-borrow-and-crowding:1}
A stock has 80 million shares, 10\% of them sold short, and trades 0.8 million shares a day. What are its days to cover?
\end{exercise}

\begin{exercise}[$\star$]\label{exo:s1:short-interest-borrow-and-crowding:2}
How can short interest exceed 100\% of the float? Give a two-step example.
\end{exercise}

\begin{exercise}[$\star$]\label{exo:s1:short-interest-borrow-and-crowding:3}
With the chapter's fee schedule, what does a stock at 80\% utilisation cost to borrow, and what does a short position of 5\% of capital in it cost for
nine months?
\end{exercise}

\begin{exercise}[$\star\star$]\label{exo:s1:short-interest-borrow-and-crowding:4}
Why does the fee book lose while the short-interest book wins, when both point at the same names?
\end{exercise}

\begin{exercise}[$\star\star$]\label{exo:s1:short-interest-borrow-and-crowding:5}
Why is the avoidance version of the signal more robust than the long--short version?
\end{exercise}

\begin{exercise}[$\star\star$]\label{exo:s1:short-interest-borrow-and-crowding:6}
A book is short twenty crowded names equally with half its capital. One of them rises twentyfold. What does the book lose?
\end{exercise}

\begin{exercise}[$\star\star\star$]\label{exo:s1:short-interest-borrow-and-crowding:7}
\emph{Coding.} Run \texttt{squeeze(0.25, 5)} and \texttt{squeeze(0.5, 10)}. Explain how the losses depend on the covered share and the speed of covering
under the square-root law.
\end{exercise}

\begin{exercise}[$\star\star\star$]\label{exo:s1:short-interest-borrow-and-crowding:8}
\emph{Find the flaw.} ``Our short-interest backtest uses FINRA short interest on its settlement date and earns 1.5\% a month.''
\end{exercise}

\section{Problem: January 2021}

\begin{problem}[{Weekend problem --- the signal and the squeeze}]\label{pb:s1:short-interest-borrow-and-crowding:1}
The synthetic lending market and the public record.

\medskip\noindent\textbf{Part I --- The episode.}
\begin{enumerate}
\item What did GameStop's short interest reach, and how can it exceed 100\%?
\item Describe the price path from January 11 to 28, 2021.
\item What did the SEC's staff conclude about short covering?
\item How is short interest reported, and how stale is it?
\end{enumerate}

\medskip\noindent\textbf{Part II --- The signal.}
\begin{enumerate}[resume]
\item What did Boehmer, Jones and Zhang, and Asquith, Pathak and Ritter, find?
\item Describe the synthetic lending market.
\item Give the ICs of short interest, days to cover, the fee and the crowding score.
\item Give the short-interest book's returns and Sharpe ratios before and after fees.
\end{enumerate}

\medskip\noindent\textbf{Part III --- The costs.}
\begin{enumerate}[resume]
\item Why does the fee book lose?
\item What does avoidance earn?
\item What did Engelberg, Reed and Ringgenberg find about short-selling risk?
\item How are recalls modelled, and why do they concentrate in the signal's names?
\end{enumerate}

\medskip\noindent\textbf{Part IV --- The verdict.}
\begin{enumerate}[resume]
\item State the \emph{named result}: the return of the high-short-interest signal before and after borrow fees, and the loss of a crowded short book in a
simulated squeeze.
\item Why are the squeeze losses floors, not forecasts?
\item What does one twentyfold name do to a crowded book?
\item How would you size a short book's crowded names?
\item Which strategy file would a long-only manager use?
\item What data would you buy to run these strategies?
\item What does the GameStop episode teach about tail risk in backtests?
\item In one sentence: what does a short seller pay for?
\end{enumerate}
\end{problem}

\section{Interview questions}

\begin{interviewq}[$\star$ \iqroles{researcher, trader}]\label{iq:s1:short-interest-borrow-and-crowding:1}
What does high short interest tell you about a stock?
\end{interviewq}

\begin{interviewq}[$\star\star$ \iqroles{trader, risk}]\label{iq:s1:short-interest-borrow-and-crowding:2}
What risks does a short position carry that a long position does not?
\end{interviewq}

\begin{interviewq}[$\star\star$ \iqroles{risk}]\label{iq:s1:short-interest-borrow-and-crowding:3}
How would you measure and limit the crowding of a short book?
\end{interviewq}

\begin{interviewq}[$\star\star$ \iqroles{researcher}]\label{iq:s1:short-interest-borrow-and-crowding:4}
Why might a signal that predicts underperformance still lose money when traded on the short side?
\end{interviewq}

\begin{interviewq}[$\star\star$ \iqroles{developer, researcher}]\label{iq:s1:short-interest-borrow-and-crowding:5}
What timestamps does a short-interest backtest need, and what goes wrong without them?
\end{interviewq}

\begin{interviewq}[$\star\star\star$ \iqroles{researcher, risk}]\label{iq:s1:short-interest-borrow-and-crowding:6}
Short sellers hold $Q$ shares of a stock trading $V$ a day with daily volatility $\sigma$. Under square-root impact, how does the price move if they all
cover within $d$ days, and why does the model understate a squeeze?
\end{interviewq}
